\documentclass[final,xcolor=table]{beamer}
\mode<presentation> {
	\usetheme{wtsi}
}
\usepackage[utf8]{inputenc}
\usepackage[T1]{fontenc}
\usepackage{lmodern} % load a font with all the characters
\usepackage[english]{babel}
\usepackage{amsmath,amsthm, amssymb, latexsym}
\usepackage{tikz}
\usetikzlibrary{shapes.geometric, arrows}
\usepackage[orientation=portrait,size=a0,scale=1.4,debug]{beamerposter}
\usefonttheme[onlymath]{serif}
\boldmath
\setbeamertemplate{caption}[numbered]

\title{30x vs 15x - What is the difference?}
\author{Martin O. Pollard, Joshua C. Randall, Manjinder S. Sandhu, Inês Barroso, Eleftheria Zeggini, George Dedoussis, Jeffrey C. Barrett}
\institute[Wellcome Trust Sanger Institute]{Wellcome Trust Sanger Institute}
\date{Today}

\begin{document}
\newcommand{\usericon}[1]{\includegraphics[width=#1\textwidth]{../images/person}}
\begin{frame}{}
    \begin{columns}[t]
    % ---------------------------------------------------------%
    % Set up a column
    \begin{column}{.49\textwidth}
        \begin{beamercolorbox}[center,wd=\textwidth]{postercolumn}
            \begin{minipage}[T]{.95\textwidth}  % tweaks the width, makes a new \textwidth
            \begin{block}{Introduction}
                Whilst the cost of sequencing has decreased substantially in recent years, cost and rate limitations remain a non-negligible factor and thus care must be taken both during sample selection and with the choice of experiment parameters, such as intended coverage, to take into account these limitations. By so doing we ensure that our experiments yield the strongest possible evidence for our experimental hypotheses from a given amount of funding.

            \end{block}
            \begin{block}{Research Question}
\begin{figure}
\centering
\begin{tikzpicture}[node distance=5.5cm]
\node (sci1) {\usericon{0.1}};
\node (lab2)[below of=sci1, yshift=-2.5cm] {High Coverage};
\node (sci2a)[right of=sci1, xshift=5cm, yshift=3cm] {\usericon{0.05}};
\node (sci2b)[below of=sci2a] {\usericon{0.05}};
\node (lab2)[below of=sci2b] {Medium Coverage};
\node (sci4a)[right of=sci2a, xshift=3.5cm] {\usericon{0.025}};
\node (sci4b)[below of=sci4a] {\usericon{0.025}};
\node (sci4c)[right of=sci4a] {\usericon{0.025}};
\node (sci4d)[below of=sci4c] {\usericon{0.025}};
\node (lab2)[below of=sci4d, xshift=-2.6cm] {Low Coverage};
\end{tikzpicture}
\end{figure}
                Can we use a reduce per-sample depth of coverage to increase total sample size without losing excessive per sample information content?
            \end{block}
            \begin{block}{Method}
                \begin{itemize} \itemindent80pt
                    \item 100 samples were taken from general Greek population
                    \item Sequenced to 30x+ on Illumina HiSeq X
                    \item Mapped using BWA MEM 0.7.12 to hs37d5 by Sanger Core Sequencing
                    \item Downsampled to 30x and 15x using samtools 1.2
                    \item Variant Calling and Genotyped Using GATK 3.x
                    \item Optional filtering with VQSR
                \end{itemize}
                {The effects of a reduction in coverage was assessed by comparing taking the 30x data for each sample as a truth set and comparing it with the 15x data for the same sample. The comparison was done using RTG vcfeval to group the variants into TP, FP, and FN and then using bcftools stats to count.}
            \end{block}
            \begin{block}{Results}
                \begin{figure}
                \includegraphics[width=0.95\textwidth]{../data/downsampling}
                \caption{TPR for NA12878 based on GIAB 0.2 at various depths of sequencing}
                \label{fig:na12878_ds}
                \end{figure}
{In our previous experiments[1] shown in figure \ref{fig:na12878_ds} 15x was shown to deliver a 95\% of the SNVs in GAIB whilst 30x delivered 98\% of the SNVs. For a convincing result however, a larger sample size is needed.}
                \begin{figure}
                \includegraphics[width=0.47\textwidth]{../data/all_present_snp}
                \includegraphics[width=0.47\textwidth]{../data/all_present_indel}
                \caption{Percentage of variants present in both 15x and 30x prior to VQSR}
                \label{fig:all_present}
                \end{figure}
{Figures \ref{fig:all_present} and \ref{fig:vqsr_present} show that the vast majority of SNVs captured in the 30x set were also captured in the 15x set. Similarly a large proportion of the}
            \end{block}
            \vfill
          % ---------------------------------------------------------%
          % end the column

            \end{minipage}
        \end{beamercolorbox}
    \end{column}
    \begin{column}{.49\textwidth}
        \begin{beamercolorbox}[center,wd=\textwidth]{postercolumn}
            \begin{minipage}[T]{.95\textwidth}  % tweaks the width, makes a new \textwidth

            \begin{block}{Results (continued)}
                \begin{figure}
                \includegraphics[width=0.47\textwidth]{../data/vqsr_present_snp}
                \includegraphics[width=0.47\textwidth]{../data/vqsr_present_indel}
                \caption{Percentage of variants present in both 15x and 30x after VQSR}
                \label{fig:vqsr_present}
                \end{figure}

                {INDEL sites were present in both sets, however the presence of a significant number of sites called variant in the 15x but not 30x INDEL results suggest that there is quite a bit of artefactual noise. This is likely down to the fact that the Illumina HiSeq is an amplification type sequencing system and also employs PCR in its library preparation which is subject to slippage.}

                {In figure \ref{fig:vqsr_present} these noise sites are still present indicating that VQSR is ineffective at identifying and removing these false positives.}

                \begin{figure}
                \includegraphics[width=0.95\textwidth]{../data/vqsr_r2}
                \caption{Dosage $R^2$ of Variants present in both 15x and 30x post VQSR}
                \label{fig:r2}
                \end{figure}

                Figure \ref{fig:r2} shows that where variants were captured in both sets, there was an $R^2$ correlation of 0.99 between the genotypes. INDELs were less promising, with a greater difference in the sites captured between the two sets and an $R^2$ correlation of 0.90 for sites present in both. There are still accuracy gains to be had from sequencing for INDELs.

            \end{block}
            \begin{block}{Conclusions}
                Sequencing at 15x coverage provides results that are good enough for population genetics worthwhile for nearly 2x the samples. SNVs don't show much improvement in accuracy above 15x. INDEL VQSR is unreliable and needs further filtering steps subsequent to that. Given observed PCR issues, PCR-free methods need to become the norm.
            \end{block}
            \begin{block}{Acknowledgements}
                Inês, Jeff and Manj for advice and guidance. Ele for samples. Joshua Randall for R and ggplot2 support
            \end{block}
            \begin{block}{References}
            %    [1] Zook, Justin M., Brad Chapman, Jason Wang, David Mittelman, Oliver Hofmann, Winston Hide, and Marc Salit. "Integrating Human Sequence Data Sets Provides a Resource of Benchmark Snp and Indel Genotype Calls." Nat Biotech 32, no. 3 (03//print 2014): 246-51.
                [1] M. O. Pollard, T. M. Keane, S. A. McCarthy, J. C. Randall, R. M. Durbin "Using haploid human DNA to design and evaluate the HiSeq X data processing strategy." ASHG 2014 Abstracts (1600M)

            \end{block}
            \vfill

          % ---------------------------------------------------------%
          % end the column

            \end{minipage}
        \end{beamercolorbox}
    \end{column}
    \end{columns}

\end{frame}
\end{document}
